% =========================================================================
% Audience cheat-sheet -- Signal Processing for Beyond-5G and 6G
% 2-page A4, article class, two columns, compact.
% Build:  pdflatex -interaction=nonstopmode cheatsheet.tex   (run twice)
% =========================================================================
\documentclass[a4paper,9pt]{extarticle}

\usepackage[a4paper,margin=10mm,columnsep=6mm]{geometry}
\usepackage[T1]{fontenc}
\usepackage{lmodern}
\usepackage{amsmath,amssymb,bm}
\usepackage{booktabs}
\usepackage{multicol}
\usepackage{xcolor}
\usepackage{enumitem}
\usepackage{titlesec}
\usepackage[hidelinks]{hyperref}

% ---- palette (matches the deck) -----------------------------------------
\definecolor{accent}{HTML}{0F7173}
\definecolor{cblue}{HTML}{2A78D6}
\definecolor{corange}{HTML}{EB6834}

% ---- compact spacing ----------------------------------------------------
\setlength{\parindent}{0pt}
\setlength{\parskip}{0pt}
\setlength{\abovedisplayskip}{2pt}
\setlength{\belowdisplayskip}{2pt}
\setlength{\abovedisplayshortskip}{1pt}
\setlength{\belowdisplayshortskip}{1pt}
\setlength{\jot}{1.5pt}
\setlist{nosep,leftmargin=1.1em}

% ---- section headers ----------------------------------------------------
\titleformat{\section}{\normalfont\bfseries\color{accent}\small}{}{0pt}{}
\titlespacing{\section}{0pt}{4pt}{1pt}

\newcommand{\thm}[1]{\vspace{1pt}\section*{}\textbf{\color{accent}#1}\par\nobreak\vspace{1pt}}
% theme header inside multicols
\newcommand{\block}[1]{\smallskip{\color{accent}\bfseries #1}\par\nobreak\vspace{1pt}}

\pagestyle{empty}

\begin{document}
\small

% ======================= TITLE ==========================================
\begin{center}
  {\large\bfseries Signal Processing for Beyond-5G / 6G --- Cheat Sheet}\\[1pt]
  {\footnotesize\color{accent} Waveforms $\bullet$ ultra-massive MIMO $\bullet$
   RIS $\bullet$ ISAC $\bullet$ AI-native PHY \quad|\quad companion to the
   90-min invited talk}\\[1pt]
  {\footnotesize Dr.\ Markkandan~S, Associate Professor, School of Electronics
   Engineering, VIT Chennai \ $\bullet$\ \texttt{markkandan.s@vit.ac.in}}
\end{center}
\vspace{1pt}
\hrule
\vspace{3pt}

\begin{multicols}{2}

% ---------------- SPECTRUM & HARDWARE -----------------------------------
\block{1. Sub-THz channel \& hardware (100--300 GHz)}
Path loss = free-space spreading $\times$ molecular absorption:
\begin{equation*}
  \mathrm{PL}(f,d)=
  \underbrace{\Big(\tfrac{4\pi f d}{c}\Big)^{2}}_{\text{spreading}}\cdot
  \underbrace{e^{\kappa_{\mathrm{abs}}(f)\,d}}_{\text{molecular abs.}}
\end{equation*}
O$_2$/H$_2$O lines at 60, 119, 183, 325 GHz carve transmission windows;
channel sparse, LoS-dominated $\Rightarrow$ pencil beams.

\smallskip
Rapp PA (AM/AM):
\begin{equation*}
  g(A)=\frac{G\,A}{\Big[1+\big(\tfrac{GA}{A_{\mathrm{sat}}}\big)^{2p}\Big]^{1/2p}}
\end{equation*}
Wiener phase noise:
\begin{equation*}
  \varphi[n{+}1]=\varphi[n]+w[n],\quad
  w\sim\mathcal N\!\big(0,\,4\pi^{2}f_c^{2}c\,T_s\big)
\end{equation*}
PA back-off $\Rightarrow$ PAPR is first-class; phase-noise power $\propto f_c^{2}$.

% ---------------- WAVEFORMS ---------------------------------------------
\block{2. Waveforms}
CP-OFDM symbol:
\begin{equation*}
  x[n]=\tfrac{1}{\sqrt N}\sum_{k=0}^{N-1}X_k\,e^{j2\pi kn/N}
\end{equation*}
OFDM ICI power under Doppler spread $f_d$ (symbol $T_s$):
\begin{equation*}
  P_{\mathrm{ICI}}\;\lesssim\;\frac{(\pi f_d T_s)^{2}}{3}
\end{equation*}
DFT-s-OFDM: $\mathbf X=F_M\mathbf s$ ($\approx$3 dB lower PAPR).

\smallskip
\textit{OTFS} --- doubly-dispersive input/output:
\begin{equation*}
  r(t)=\iint h(\tau,\nu)\,s(t-\tau)\,e^{j2\pi\nu(t-\tau)}\,\mathrm d\tau\,\mathrm d\nu
\end{equation*}
ISFFT (DD $\to$ TF grid):
\begin{equation*}
  X_{\mathrm{tf}}[n,m]=\tfrac{1}{\sqrt{NM}}
  \sum_{k=0}^{N-1}\sum_{l=0}^{M-1}\! x[k,l]\,
  e^{\,j2\pi(\frac{nk}{N}-\frac{ml}{M})}
\end{equation*}
DD input--output (2D twisted convolution):
\begin{equation*}
  y[k,l]\approx\sum_i h_i\,x\big[(k{-}k_i)_N,(l{-}l_i)_M\big]e^{j2\pi\phi_i}+w[k,l]
\end{equation*}
SFFT is the inverse of ISFFT; channel is sparse \& quasi-static in DD.

\smallskip
\textit{AFDM} --- discrete affine Fourier (chirp) basis:
\begin{equation*}
  \phi_m[n]=\tfrac{1}{\sqrt N}\,e^{\,j2\pi(c_1 n^{2}+\frac{mn}{N}+c_2 m^{2})}
\end{equation*}
Tune $c_1$ to the Doppler span; paths separate in the DAFT domain with
1D single-symbol processing.

% ---------------- MIMO --------------------------------------------------
\block{3. Ultra-massive MIMO \& near field}
Rayleigh (Fraunhofer) distance, aperture $D$:
\begin{equation*}
  d_R=\frac{2D^{2}}{\lambda}\qquad
  \begin{cases} d>d_R:&\text{planar }\mathbf a(\theta)\\
                d<d_R:&\text{spherical }\mathbf a(\theta,r)\end{cases}
\end{equation*}
$D{=}0.5$ m @ 30 GHz $\Rightarrow d_R{=}50$ m: users \emph{inside} near field
$\Rightarrow$ range-selective beam focusing.

\smallskip
Hybrid precoding (unit-modulus RF, $|[F_{\mathrm{RF}}]_{ij}|{=}1/\sqrt{N_t}$):
\begin{equation*}
  \min_{F_{\mathrm{RF}},F_{\mathrm{BB}}}
  \big\|F_{\mathrm{opt}}-F_{\mathrm{RF}}F_{\mathrm{BB}}\big\|_F
  \;\text{s.t.}\; F_{\mathrm{RF}}\!\in\!\{\text{array responses}\}
\end{equation*}
Angle-sparse mmWave $\Rightarrow$ OMP; $N_{\mathrm{RF}}\!\approx\!N_s$
chains $\approx$ digital.

\smallskip
CS channel estimation \& pilot contamination:
\begin{equation*}
  \mathbf y=\Phi\mathbf h+\mathbf n,\;\|\mathbf h\|_0{=}L\!\ll\!\dim,\;\;
  \mathrm{SINR}_j\!\xrightarrow{M\to\infty}\!\frac{\beta_{jj}^{2}}{\sum_{l\neq j}\beta_{jl}^{2}}
\end{equation*}

% ---------------- RIS ---------------------------------------------------
\block{4. Reconfigurable intelligent surfaces}
Cascaded channel ($N$ elements):
\begin{equation*}
  y=\big(\bm h_r^{H}\Theta\,\bm G+h_d^{H}\big)x+n,\;\;
  \Theta=\mathrm{diag}(\beta_1 e^{j\theta_1},\dots,\beta_N e^{j\theta_N})
\end{equation*}
Optimal phases $\Rightarrow$ received power scales as $N^{2}$:
\begin{equation*}
  \theta_n^{\star}=\arg(h_d)-\arg\!\big([\bm h_r]_n^{*}[\bm G]_n\big)
  \;\Rightarrow\;
  P_{\mathrm{rx}}\!\propto\!\Big(\textstyle\sum_n|[\bm h_r]_n||[\bm G]_n|\Big)^{2}\!\!\sim\! N^{2}
\end{equation*}
Co-phasing: $+6$ dB per doubling; random phases $\Rightarrow P\propto N$.
Product path loss $\propto(d_1 d_2)^{-\alpha}$ (passivity tax).

% ---------------- ISAC --------------------------------------------------
\block{5. Radar \& ISAC}
Delay / Doppler / resolutions:
\begin{equation*}
  \tau=\frac{2R}{c},\;\; f_D=\frac{2v f_c}{c},\;\;
  \Delta R=\frac{c}{2B},\;\; \Delta v=\frac{c}{2 f_c T_{\mathrm{obs}}}
\end{equation*}
OFDM radar (RX knows $X$): $D[n,m]=Y[n,m]/X[n,m]$, then IFFT over $n$ /
FFT over $m$ $\to$ range--Doppler map. Delay CRB:
\begin{equation*}
  \mathrm{var}(\hat\tau)\;\ge\;
  \frac{1}{8\pi^{2}\,\beta_{\mathrm{rms}}^{2}\,\mathrm{SNR}\cdot N_{\mathrm{obs}}}
\end{equation*}
($\beta_{\mathrm{rms}}$ = RMS bandwidth $\Rightarrow$ edge subcarriers are gold.)

% ---------------- ACCESS ------------------------------------------------
\block{6. Multiple access}
NOMA (SIC, superposition; near user cancels far):
\begin{align*}
  R_{\mathrm{far}}&=\log_2\!\Big(1+\tfrac{P_2|h_f|^2}{P_1|h_f|^2+\sigma^2}\Big),\\
  R_{\mathrm{near}}&=\log_2\!\Big(1+\tfrac{P_1|h_n|^2}{\sigma^2}\Big)
\end{align*}
RSMA (split into common + private):
\begin{equation*}
  R_k=C_k+R_{p,k},\quad
  \textstyle\sum_k C_k\le\min_k\log_2(1+\mathrm{SINR}_{c,k})
\end{equation*}

% ---------------- NTN & EE ----------------------------------------------
\block{7. NTN \& energy efficiency}
LEO Doppler ($v_{\mathrm{sat}}\!\approx\!7.6$ km/s):
\begin{equation*}
  f_D(t)=\frac{f_c}{c}\,v_{\mathrm{rel}}(t)\;\Rightarrow\;\pm25\,\text{ppm}
  \;(\pm50\,\text{kHz @ 2 GHz})
\end{equation*}
Energy efficiency \& ADC power:
\begin{equation*}
  \eta_{EE}=\frac{R}{P_{\mathrm{tot}}}\ \text{[bit/J]},\qquad
  P_{\mathrm{ADC}}\propto 2^{b} f_s
\end{equation*}

\end{multicols}

\vspace{1pt}
\hrule
\vspace{3pt}

% ---------------- TECH COMPARISON TABLE ---------------------------------
{\color{accent}\bfseries Waveform comparison at a glance}\quad
{\footnotesize (relative, qualitative)}
\vspace{2pt}

\begin{center}
\renewcommand{\arraystretch}{1.15}
\footnotesize
\begin{tabular}{@{}lcccc@{}}
\toprule
\textbf{} & \textbf{CP-OFDM} & \textbf{DFT-s-OFDM} & \textbf{OTFS} & \textbf{AFDM} \\
\midrule
PAPR                 & high        & \textbf{low} ($-$3 dB) & high        & moderate \\
Doppler robustness   & poor        & poor                   & \textbf{excellent} & \textbf{excellent} \\
Diversity (doubly-disp.) & low     & low                    & full        & full \\
Rx / detection cost  & \textbf{low} (1-tap) & low + FDE       & high (2D MP) & moderate (1D) \\
Processing domain    & TF          & TF (SC)                & delay--Doppler & DAFT (chirp) \\
Pilot / CE overhead  & moderate    & moderate               & low (DD impulse) & low \\
Maturity / standards & \textbf{deployed} (4G/5G) & \textbf{5G UL} & research     & emerging (2023$+$) \\
\bottomrule
\end{tabular}
\end{center}
\vspace{1pt}
{\footnotesize No single waveform wins every column $\Rightarrow$ expect
scenario-adaptive PHY. Three formulas organize half the field:\;
$P\!\propto\!N^{2}$ \ $\bullet$\ $d_R=2D^{2}/\lambda$ \ $\bullet$\
$\mathrm{var}(\hat\tau)\ge 1/(8\pi^{2}\beta^{2}\,\mathrm{SNR})$.}

\vspace{3pt}
\hrule
\vspace{3pt}

% ---------------- AI-NATIVE + KEY NUMBERS + DEMOS ------------------------
\begin{multicols}{2}

\block{8. AI-native physical layer}
\emph{Model deficit} (no faithful model, e.g.\ PA nonlinearity) vs
\emph{algorithm deficit} (optimum intractable) $\Rightarrow$ learn the model /
the algorithm. \textbf{Deep unfolding} turns $K$ iterations into a $K$-layer net:
\begin{align*}
  \text{ISTA:}\;& \mathbf x^{(k+1)}=\eta_{\lambda}\big(\mathbf x^{(k)}+A^{H}(\mathbf y-A\mathbf x^{(k)})\big)\\
  \text{LISTA:}\;& \mathbf x^{(k+1)}=\eta_{\lambda_k}\big(W_1^{(k)}\mathbf y+W_2^{(k)}\mathbf x^{(k)}\big)
\end{align*}
Learned CSI estimation $\approx$ MMSE \emph{without} knowing statistics; neural
receivers run inside 3GPP frames; end-to-end autoencoder PHY learns the
constellation. 3GPP Rel-18/19: CSI-feedback compression, beam prediction,
positioning. Open: data, generalization, trust.

\block{Key numbers to remember}
\begin{itemize}[leftmargin=1.1em,itemsep=0.3pt]
  \item 500 km/h @ 28 GHz $\Rightarrow f_d\!\approx\!13$ kHz ($\varepsilon\!\approx\!0.11$)
  \item $D{=}0.5$ m @ 30 GHz $\Rightarrow d_R{=}50$ m (near field)
  \item $B{=}30$ MHz $\Rightarrow \Delta R\!\approx\!5$ m
  \item RIS $N^{2}$: $+6$ dB per doubling of $N$
  \item LEO 600 km @ 2 GHz $\Rightarrow \pm50$ kHz swing
  \item DFT-s-OFDM $\approx3$ dB less PAPR than CP-OFDM
  \item 3GPP: Rel-19 (ISAC model, AI/ML) $\to$ Rel-20/21 (6G) $\to$ spec $\sim$2029
\end{itemize}

\block{Five live MATLAB demos}
\begin{enumerate}[leftmargin=1.4em,itemsep=0.3pt]
  \item OFDM vs OTFS @ 500 km/h: OFDM BER floor, OTFS keeps falling
  \item 64-el hybrid BF: 4 RF chains $\approx$ fully digital
  \item RIS $N{=}16{-}1024$: slopes 2 (opt.) vs 1 (random)
  \item OFDM-ISAC: two moving targets on a range--Doppler map
  \item LS / MMSE / NN channel est.: learned $\approx$ MMSE
\end{enumerate}

\end{multicols}

\vspace{2pt}
\hrule
\vspace{3pt}

% ---------------- READING LIST ------------------------------------------
{\color{accent}\bfseries Reading list --- 15 landmark references \& surveys}
\vspace{2pt}

\begin{multicols}{2}
\footnotesize
\begin{enumerate}[leftmargin=1.4em,itemsep=0.5pt]
  \item ITU-R Rec.\ M.2160-0, ``Framework and overall objectives of the
        future development of IMT for 2030 and beyond,'' Nov.\ 2023.
  \item W. Saad, M. Bennis, and M. Chen, ``A vision of 6G wireless
        systems: Applications, trends, technologies, and open research
        problems,'' \emph{IEEE Network}, vol.\ 34, no.\ 3, 2020.
  \item H. Tataria \emph{et al.}, ``6G wireless systems: Vision,
        requirements, challenges, insights, and opportunities,''
        \emph{Proc.\ IEEE}, vol.\ 109, no.\ 7, 2021.
  \item R. Hadani \emph{et al.}, ``Orthogonal time frequency space
        modulation,'' in \emph{Proc.\ IEEE WCNC}, 2017.
  \item Z. Wei \emph{et al.}, ``Orthogonal time-frequency space
        modulation: A promising next-generation waveform,'' \emph{IEEE
        Wireless Commun.}, vol.\ 28, no.\ 4, 2021.
  \item A. Bemani, N. Ksairi, and M. Kountouris, ``AFDM: Affine frequency
        division multiplexing for next generation wireless
        communications,'' \emph{IEEE Trans.\ Wireless Commun.}, 2023.
  \item O. El Ayach, S. Rajagopal, S. Abu-Surra, Z. Pi, and R. W. Heath,
        ``Spatially sparse precoding in millimeter wave MIMO systems,''
        \emph{IEEE Trans.\ Wireless Commun.}, vol.\ 13, no.\ 3, 2014.
  \item E. Bj\"ornson, J. Hoydis, and L. Sanguinetti, ``Massive MIMO has
        unlimited capacity,'' \emph{IEEE Trans.\ Wireless Commun.},
        vol.\ 17, no.\ 1, 2018.
  \item Q. Wu and R. Zhang, ``Intelligent reflecting surface enhanced
        wireless network via joint active and passive beamforming,''
        \emph{IEEE Trans.\ Wireless Commun.}, vol.\ 18, no.\ 11, 2019.
  \item M. Di Renzo \emph{et al.}, ``Smart radio environments empowered
        by reconfigurable intelligent surfaces,'' \emph{IEEE J.\ Sel.\
        Areas Commun.}, vol.\ 38, no.\ 11, 2020.
  \item F. Liu \emph{et al.}, ``Integrated sensing and communications:
        Toward dual-functional wireless networks for 6G and beyond,''
        \emph{IEEE J.\ Sel.\ Areas Commun.}, vol.\ 40, no.\ 6, 2022.
  \item C. Sturm and W. Wiesbeck, ``Waveform design and signal processing
        aspects for fusion of wireless communications and radar
        sensing,'' \emph{Proc.\ IEEE}, vol.\ 99, no.\ 7, 2011.
  \item T. O'Shea and J. Hoydis, ``An introduction to deep learning for
        the physical layer,'' \emph{IEEE Trans.\ Cogn.\ Commun.\ Netw.},
        vol.\ 3, no.\ 4, 2017.
  \item N. Shlezinger, J. Whang, Y. C. Eldar, and A. J. Goldsmith,
        ``Model-based deep learning,'' \emph{Proc.\ IEEE}, vol.\ 111,
        no.\ 5, 2023.
  \item Y. Mao, O. Dizdar, B. Clerckx, R. Schober, P. Popovski, and
        H. V. Poor, ``Rate-splitting multiple access: Fundamentals,
        survey, and future research trends,'' \emph{IEEE Commun.\
        Surveys Tuts.}, vol.\ 24, no.\ 4, 2022.
\end{enumerate}
\end{multicols}

\vspace{1pt}
{\footnotesize\color{gray}Tools: MATLAB (SPT/CommT/DLT) $\bullet$ NVIDIA Sionna
$\bullet$ this talk's code \& Python verification ship together.}

\end{document}
